\section{Introduction}

Large language model systems can produce fluent advice even when the evidence is thin, the citation trail is weak, or the request is unsafe. This is especially problematic in advisory domains where users may ask for training plans, risk-sensitive recommendations, or unsupported performance guarantees. Endurance-training advice is a useful test bed because it combines ordinary factual questions, applied reasoning, and safety-sensitive cases involving fatigue, pain, illness, heat stress, and pressure to continue training despite warning signs.

Ordinary RAG evaluation is not enough for this setting. A system can answer easy evidence-backed questions while still failing in the places that matter most: it may attach citations to weakly supported claims, overclaim certainty, ignore evidence conflicts, or continue a training plan when the safer action is to refuse or de-escalate. In such a domain, refusal is not merely an error mode. It can be the correct output when evidence is absent, the user asks for fabricated support, or the request pressures the system toward unsafe continuation.

We therefore study a narrower workflow question: can an evidence-gated audit-and-repair pipeline make grounding, refusal, repair, and safety-boundary behavior more observable and partially controllable? Our system decomposes advisory RAG into explicit control points: request-level filtering, evidence retrieval, evidence sufficiency checking, risk checking, evidence-constrained generation, independent grounding audit, and bounded repair or refusal. These stages are logged as trace fields, so failures are not only visible in the final answer but also in the intermediate decisions that produced it.

The contribution is not a new foundation model, a general multi-agent framework, or a universal safety defense. It is a workflow-level diagnostic structure for safety-sensitive advisory RAG. The system treats \texttt{answered}, \texttt{partial\_answer}, and \texttt{refused} as separate output states, allowing citation repair, conservative false refusal, designed-unanswerable refusal, and request-level stress blocking to be measured rather than hidden inside a single success/failure score.

Our main contributions are fourfold. First, we propose an evidence-gated audit-and-repair workflow for safety-sensitive advisory RAG. Second, we introduce a 100-question pilot benchmark for endurance-training advice, with 90 evidence-backed questions and 10 designed-unanswerable controls. Third, we build a 30-prompt targeted safety-stress suite covering prompt injection, unsafe requests, citation hallucination pressure, overclaim requests, and evidence-conflict pressure. Fourth, we report diagnostic runs, ablations, and case studies that show how refusal, repair, and targeted hardening behave in practice.

We keep the scope deliberately narrow. The benchmark is pilot-scale and author-authored, the primary result uses one local model with a supplementary llama3 subset and stress check, and the deterministic hardening layer targets known request patterns rather than general adversarial prompts. The results should therefore be read as evidence for a traceable workflow-level diagnostic approach, not as clinical validation, coaching efficacy, deployment readiness, or broad prompt-injection robustness.
